\subsection{\foreign{Boolean Satisfiability}}
\label{subsec:bs}

Pertama-tama, didefinisikan pemetaan injektif \(\boolmap\) dari \(V\) ke variabel boolean di SMT. Catat bahwa variabel boolean di sini bukanlah elemen di \(\Bbb\), melainkan merupakan variabel abstrak di SMT yang nilainya belum ditentukan. Hanya setelah menjalankan \foreign{solver} barulah dihasilkan \foreign{model} yang elemennya di \(\Bbb\). Dengan demikian, variabel-variabel tersebut digunakan untuk mengkonstruksi \foreign{constraint} sehingga \foreign{model} \(\in \Bbb^{|V|}\) menunjukkan \(W \subseteq V\) yang elemen-elemennya berkoresponden dengan elemen di \(V\) yang \(\boolmap(v_i)\)-nya bernilai \(\True\).

\clearpage

Meskipun secara konsep \(\boolmap(v_i)\) adalah variabel abstrak, secara matematis, tidak dibedakan makna \(\boolmap(v_i)\) antara sebelum dan setelah adanya \foreign{model}. Oleh karena itu, secara praktis, tidak ada salahnya menganggap \(\boolmap(v_i) \in \Bbb\). Cukup dibedakan pemaknaan penggunaan ketika dalam konteks \foreign{constraint}. Menggunakan konvensi ini, \(\boolmap(V) = X\) dapat digunakan langsung untuk \foreign{indexing} pada \(V\), sehingga \(W = V[X] \subseteq V\) terdefinisi dengan baik. Sebelum ada \foreign{model}, \(W\) adalah koresponden dari ruang solusi terhadap \foreign{constraint}, dan setelah ada \foreign{model}, solusi \(W\) adalah himpunan nyata.

\begin{example}
	Bila \(\boolmap(v_i) = x_i\), maka untuk simpul-simpul di \autoref{fig:algorithm-example}, didapatkan \(\boolmap(V) = \{x_1, x_2, x_3, x_4, x_5\} = X\). Kemudian, model \(x_3 = x_4 = x_5 = \True, x_1 = x_2 = \False\), menunjukkan koresponden himpunan \(W = \{x_3, x_4, x_5\} \subseteq V\). Dapat digunakan notasi alternatif \(X = (\False, \False, \True, \True, \True)\) menunjukkan model dan \(W = V[X]\) sebagai koresponden himpunannya.
\end{example}

Kedua, untuk sembarang \(W \subseteq V\), dapat dikonstruksi \foreign{constraint} sedemikian rupa \foreign{model}-nya berkoresponden satu-satu dengan \(W\). Secara umum, variabel yang tidak memiliki restriksi akan memiliki nilai \(\indeterminate\) pada \foreign{model}-nya. Maksud dari \(\indeterminate\) adalah variabel yang bebas; yaitu, tidak peduli apa pun nilai yang ia miliki, kebenaran dari \foreign{constraint} yang ada tidak terpengaruhi oleh variabel tersebut. Ketika \foreign{solver} menghasilkan sebuah \foreign{model}, seluruh variabelnya yang \(\indeterminate\) akan diberi sembarang nilai. Artinya, untuk variabel boolean, ia akan mendapatkan sembarang nilai \(\True\) atau \(\False\). Maka dari itu, agar koresponden \foreign{model}-nya satu-satu dengan himpunan yang dituju, tidak boleh ada variabel yang tidak terestriksi secara tidak sengaja dalam \foreign{constraint} yang dikonstruksikan. Untuk kasus di sini yang memilih tepat \(W\), \foreign{constraint} yang menyatakan pemilihan ini adalah sebagai berikut:
\[
	\combine{\left(}
			{\bigwedge \limits_{\displaystyle v \in W} \boolmap(v)}
	{\right)}{\wedge}{\left(}
			{\bigwedge \limits_{\displaystyle v \in \widetilde{W}} \neg \boolmap(v)}
	{\right)}
\]

\begin{example}
	Untuk \(V = \{v_1, v_2, v_3, v_4, v_5\}\) dan \(W = \{v_3, v_4, v_5\}\), \foreign{constraint} yang menyebabkan koresponden satu-satu antara \(W\) dan \(V[X]\) adalah: \(\neg x_1 \wedge \neg x_2 \wedge x_3 \wedge x_4 \wedge x_5\).
\end{example}

Ketiga, untuk sembarang \(W \subseteq V\), dapat dikonstruksi \foreign{constraint} sedemikian rupa \foreign{model}-nya berkoresponden satu-satu dengan \(\Pcal(W)\). Perhatikan pada konstruksi \foreign{constraint} sebelumnya betapa pentingnya konteks \(V\) yang dimaksud relatif terhadap \(W\) karena \foreign{model} yang menyatakan \(W\) tidak hanya didefinisikan oleh \(W\) sendiri, tetapi juga oleh \(\widetilde{W} = V \setminus W\). Dalam konstruksi \foreign{constraint} untuk \(\Pcal(W)\) ini, hubungan tersebut menjadi lebih kritikal. Karena, terbalik dari perspektif \foreign{power set} pada umumnya, dalam bentuk \foreign{constraint}, definisi \foreign{power set} justru dilihat dari elemen-elemennya yang tidak ada di \(W\). Lebih formalnya dinyatakan oleh ekspresi berikut:
\[
	\bigwedge \limits_{\displaystyle v \in \widetilde{W}} \neg \boolmap(v)
\]
Dengan kata lain, seluruh kemungkinan \foreign{model} yang memenuhi \foreign{constraint} ini bebas menetapkan keanggotaan elemen-elemen di \(W\), dan hanya yang di \(W\), sehingga menyatakan seluruh kombinasi yang merepresentasikan \(\Pcal(W)\). Dengan demikian, ruang solusi dari \foreign{model} yang memenuhi \foreign{constraint} tersebut ekivalen dengan \(\Pcal(W)\).

\begin{example}
	Untuk \(V = \{v_1, v_2, v_3, v_4, v_5\}\) dan \(W = \{v_3, v_4, v_5\}\), \foreign{constraint} yang menyebabkan koresponden satu-satu antara \(\Pcal(W)\) dan \(V[X]\) adalah: \(\neg x_1 \wedge \neg x_2\). Setiap \foreign{model} yang memenuhi \(\neg x_1 \wedge \neg x_2\) mewakili satu subset dari \(\{v_3, v_4, v_5\}\).
\end{example}

Akhirnya, dapat dibahas konstruksi \foreign{constraint boolean satisfiability} yang menunjukkan \foreign{resolving set}. Ingat kembali, untuk setiap kolom di \(P\), himpunan \(V[P_{*j}]\) serta seluruh himpunan bagiannya bukanlah \foreign{resolving set}. Maka, \(W \subseteq V\) adalah \foreign{resolving set} jika dan hanya jika \(W\) tidak ada di \(\Pcal(V[P_{*j}])\) untuk setiap kolom di \(P\). Secara formal, ini dapat dinyatakan sebagai berikut:
\[
	\bigwedge \limits_{\displaystyle P_{*j} \in P} \neg \left( \bigwedge \limits_{\displaystyle v \in \widetilde{V[P_{*j}]}} \neg \boolmap(v) \right)
\]

Dan dapat disederhakan menjadi:
\begin{equation}
	\bigwedge \limits_{\displaystyle P_{*j} \in P} \enspace \bigvee \limits_{\displaystyle v \in V[\sim\!P_{*j}]} \boolmap(v)
	\label{eq:boolean-satisfiability-natural}
\end{equation}

\begin{example}
	Ambil graf \autoref{fig:algorithm-example}. \foreign{Constraint} yang berkorespondensi dengan \foreign{resolving set}-nya adalah:
		\[
			( x_1 \vee x_2 ) \wedge \ ( x_1 \vee x_3 ) \wedge \ ( x_2 \vee x_3 ) \wedge \ ( x_4 \vee x_5 )
			\vspace{-1.7em}
		\]
	\label{ex:boolean-satisfiability-natural}
\end{example}

Secara teori, \foreign{constraint} \autoref{eq:boolean-satisfiability-natural} sudah berbentuk \foreign{conjunctive normal form} (CNF). Akan tetapi, ketika menunjukkan \(\unsat\) menggunakan SMT, efektivitasnya sangat bergantung pada \foreign{conflict driven learning}, yaitu metode "penalaran" yang digunakan oleh \foreign{solver} untuk melakukan pemangkasan ruang pencarian. Salah satu sifat yang digunakan oleh \foreign{conflict driven learning} adalah bentuk CNF yang terdiri dari suku-suku positif maupun suku-suku negatif. Sedangkan \autoref{eq:boolean-satisfiability-natural} sepenuhnya hanya memuat suku-suku positif. Bentuk yang sudah diperoleh ini dapat ditransformasikan agar juga memuat suku-suku negatif dengan memfaktorkan keluar "pemilihan" sang himpunan kosong, dengan cara berikut:
\begin{equation}
	\begin{array}{l}
		\displaystyle \bigwedge \limits_{\displaystyle P_{*j} \in P} \enspace \combine{\left(}{\False}{\right)}{\vee}{\left(}{\bigvee \limits_{\displaystyle v \in V[\sim\!P_{*j}]} \boolmap(v)}{\right)} \\ [18pt]
		\displaystyle = \bigwedge \limits_{\displaystyle P_{*j} \in P} \enspace \combine{\left(}{\combine{\left(}{\bigvee \limits_{\displaystyle v \in V[P_{*j}]} \boolmap(v)}{\right)}{\wedge}{\left(}{\neg \bigvee \limits_{\displaystyle v \in V[P_{*j}]} \boolmap(v)}{\right)}}{\right)}{\vee}{\left(}{\bigvee \limits_{\displaystyle v \in V[\sim\!P_{*j}]} \boolmap(v)}{\right)} \\ [18pt]
		\displaystyle = \bigwedge \limits_{\displaystyle P_{*j} \in P} \enspace \combine{\left(}{\combine{\left(}{\bigvee \limits_{\displaystyle v \in V[P_{*j}]} \boolmap(v)}{\right)}{\vee}{\left(}{\bigvee \limits_{\displaystyle v \in V[\sim\!P_{*j}]} \boolmap(v)}{\right)}}{\right)}{\\[18pt]\qquad\qquad\wedge}{\left(}{\combine{\left(}{\neg \bigvee \limits_{\displaystyle v \in V[P_{*j}]} \boolmap(v)}{\right)}{\vee}{\left(}{\bigvee \limits_{\displaystyle v \in V[\sim\!P_{*j}]} \boolmap(v)}{\right)}}{\right)} \\ [18pt]
		\displaystyle = \bigwedge \limits_{\displaystyle P_{*j} \in P} \enspace \combine{\left(}{\bigvee \limits_{\displaystyle v \in V} \boolmap(v)}{\right)}{\wedge}{\left(}{\bigvee \limits_{\displaystyle v \in V[P_{*j}]} \boolmap(v)\ \implies \bigvee \limits_{\displaystyle v \in V[\sim\!P_{*j}]} \boolmap(v)}{\right)} \\ [18pt]
		\displaystyle = \combine{\left(}{\bigvee \limits_{\displaystyle v \in V} \boolmap(v)}{\right)}{\wedge}{\left(}{\bigwedge \limits_{\displaystyle P_{*j} \in P} \enspace \left( \bigvee \limits_{\displaystyle v \in V[P_{*j}]} \boolmap(v)\ \implies \bigvee \limits_{\displaystyle v \in V[\sim\!P_{*j}]} \boolmap(v) \right)}{\right)} \\ [18pt]
	\end{array}
	\label{eq:boolean-satisfiability}
\end{equation}

\begin{example}
	Berikut adalah hasil dari pemfaktoran \autoref{ex:boolean-satisfiability-natural}.
	\[
		\smash[b]{
		\begin{array}{c}
			x_1 \vee x_2 \vee x_3 \vee x_4 \vee x_5 \\ [4pt]
			x_3 \vee x_4 \vee x_5 \implies x_1 \vee x_2 \\
			x_2 \vee x_4 \vee x_5 \implies x_1 \vee x_3 \\
			x_1 \vee x_4 \vee x_5 \implies x_2 \vee x_3 \\
			x_1 \vee x_2 \vee x_3 \implies x_4 \vee x_5
		\end{array}
		} \vspace{.5em}
	\]
	\label{ex:boolean-satisfiability}
\end{example}

Karena ruang solusi \foreign{model} yang memenuhi \foreign{constraint} tersebut sudah menunjukkan \foreign{resolving set}, maka algoritma yang perlu dirancang cukup mencari \foreign{resolving set} dengan kardinalitas terendah dari ruang solusi. Ini dapat dilakukan dengan mudah menggunakan \foreign{constraint} \(\AtLeast\) dan \(\AtMost\) yang memaksa banyaknya variabel boolean bernilai \(\True\) secara berturut-turut kurang dari atau sama dengan dan lebih dari atau sama dengan \(k\).

Secara historis, SMT paling efektif ketika bekerja langsung dengan bentuk CNF. Akan tetapi, SMT \foreign{modern} sudah tidak lagi terbatasi oleh \foreign{implementation detail} seperti ini, sebagaimana yang telah dibahasa pada Pendahuluan. Maka dari itu, \foreign{solver} sekarang sudah memiliki \foreign{higher level construct} dan justru lebih banyak diuntungkan ketika bekerja langsung dalam bentuk tersebut. \foreign{Constraint} \(\AtLeast\) dan \(\AtMost\) merupakan salah satu contoh dari \foreign{higher level construct} yang ada dalam SMT. Namun, tidak ada salahnya untuk memahami bagaimana bentuk-bentuk ini dapat direpresentasikan menggunakan operasi konjungsi dan disjungsi yang lebih sederhana. Berikut adalah salah satu \foreign{encoding} yang mungkin dilakukan untuk menyatakan kedua \foreign{constraint} tersebut.
\[
	\begin{array}{r@{\ }l}
		\AtLeast(V, n) &= \displaystyle \bigvee \limits_{\displaystyle W \in \binom{V}{n}} \enspace \bigwedge \limits_{\displaystyle v \in W} \boolmap(v) \\ [48pt]
		\AtMost(V, n) &= \displaystyle \bigvee \limits_{\displaystyle W \in \binom{V}{n}} \enspace \bigwedge \limits_{\displaystyle v \in \widetilde{W}} \neg \boolmap(v)
	\end{array}
\]
Catat, notasi \(\displaystyle \binom{V}{n} = \{W \subseteq V : |V| = n\}\) menunjukkan himpunan yang memuat seluruh himpunan bagian \(V\) dengan kardinalitas \(n\).

Secara umum, menunjukkan \(\sat\) dalam SMT lebih murah dari pada menunjukkan \(\unsat\). Hal ini karena secara fundamental, SMT adalah proses mencari \foreign{model}. Oleh karena itu, untuk menunjukkan \(\sat\), cukup dicari satu saja contoh yang memenuhi \foreign{constraint}. Sedangkan untuk menunjukkan \(\unsat\), harus secara habis dibuktikan bahwa tidak ada satu pun \foreign{model} yang dapat memenuhi \foreign{constraint}.

Dengan demikian, didapatkan desain algoritma yang melakukan pencarian dari atas. Digunakan \(\AtLeast\) dan \(\AtMost\) secara bersamaan untuk memaksa mencari \foreign{model} berkardinalitas tepat sama dengan \(k\). Dimulai dari nilai \(k\) tertinggi, lalu secara iteratif menurunkan nilai \(k\) hingga menemukan kardinalitas pertama yang menyebabkan \(\unsat\), dan mengambil \(k + 1\) sebagai dimensi metrik.

\begin{example}
	Ambil graf \autoref{fig:algorithm-example}, dan cari dimensi metriknya menggunakan \autoref{alg:boolean-satisfiability}. \foreign{Constraint} yang digunakan sesuai dengan \autoref{ex:boolean-satisfiability}. Kemudian, untuk masing-masing tahap berikut, diterapkan \foreign{constraint} tambahan \(\AtLeast\) dan \(\AtMost\) sama dengan \(k\).
	\begin{enumerate}
		\item Untuk \(k = 5\), ditemukan \(\sat\).
		\item Untuk \(k = 4\), ditemukan \(\sat\).
		\item Untuk \(k = 3\), ditemukan \(\sat\).
		\item Untuk \(k = 2\), ditemukan \(\unsat\).
	\end{enumerate}
	Disimpulkan bahwa \(2 + 1 = 3\) adalah dimensi metrik dari graf tersebut.
\end{example}

\begin{algorithm}
	\foreign{Input}:
	\begin{enumerate}
		\item Representasi graf dalam representasi matriks ketetanggaan \(M\).
		\item Mencari dimensi metrik atau dimensi metrik sisi?
	\end{enumerate}
	Prosedur:
	\begin{enumerate}
		\item[1--5.] Sama seperti pada \autoref{alg:brute-force}.
		\setcounter{enumi}{5}
		\item Tambahkan \foreign{constraint} berdasarkan \autoref{eq:boolean-satisfiability} dari \(P\) ke \foreign{solver} SMT.
		\item Mulai variabel \(k\) dari banyaknya simpul, yaitu \(k := |V|\).
		\item Selama \(k \geq 0\):
			\begin{enumerate}
					\item Tambahkan \foreign{constraint} \(\AtLeast\) dan \(\AtMost\) sama dengan \(k\).
					\item Jalankan \foreign{solver} untuk mencari \foreign{model}.
					\item Jika dihasilkan \(\unsat\), lompat ke tahap (10).
					\item Hapus \foreign{constraint} \(\AtLeast\) dan \(\AtMost\) yang ditambahkan sebelumnya.
					\item Kurangi \(k\) sebanyak satu, yaitu \(k := k - 1\).
			\end{enumerate}
		\item Simpulkan bahwa 0 adalah dimensi metrik atau dimensi metrik sisi, dan \texttt{BERHENTI}.
		\item Simpulkan bahwa \(k + 1\) adalah dimensi metrik atau dimensi metrik sisi.
	\end{enumerate}
	\caption{\foreign{Boolean Satisfiability}}
	\label{alg:boolean-satisfiability}
\end{algorithm}
